\documentclass[12pt,a4paper]{article}

\usepackage[utf8]{inputenc}
\usepackage[T1]{fontenc}
\usepackage[french]{babel}
\usepackage{geometry}
\usepackage{array}
\usepackage{longtable}
\usepackage{xcolor}
\usepackage{hyperref}
\usepackage{enumitem}
\usepackage{listings}

\geometry{margin=2.5cm}

\hypersetup{
  colorlinks=true,
  linkcolor=blue,
  urlcolor=blue,
  citecolor=blue
}

\lstdefinestyle{shellstyle}{
  basicstyle=\ttfamily\small,
  breaklines=true,
  frame=single,
  columns=fullflexible,
  backgroundcolor=\color{gray!8}
}

\title{Premiers pas DevSecOps}
\author{SecureRAG Hub}
\date{\today}

\begin{document}

\maketitle

\section{Introduction}

Le DevSecOps consiste à intégrer la sécurité directement dans le cycle de
développement, d'intégration, de livraison et d'exploitation. L'objectif n'est
pas d'ajouter une étape de contrôle à la fin du projet, mais de rendre la
sécurité continue, automatisée, mesurable et vérifiable.

Dans une démarche DevSecOps moderne, chaque changement de code doit pouvoir
être construit, testé, scanné, signé, déployé et prouvé. Les preuves produites
par la chaîne CI/CD deviennent aussi importantes que le déploiement lui-même,
car elles permettent de démontrer la qualité, la traçabilité et la conformité du
système.

\section{Objectifs initiaux}

Les premiers pas DevSecOps visent à poser une base fiable avant d'aller vers des
mécanismes plus avancés comme la signature d'images, l'attestation de release,
les politiques Kubernetes ou la résilience de données.

\begin{itemize}[leftmargin=*]
  \item automatiser les tests et validations de base ;
  \item contrôler la qualité du code dès les premiers commits ;
  \item scanner les dépendances et les images conteneurisées ;
  \item standardiser les environnements de build et de déploiement ;
  \item produire des artefacts de preuve exploitables ;
  \item éviter l'exposition de secrets dans le dépôt ;
  \item préparer un déploiement reproductible et traçable.
\end{itemize}

\section{Étape 1 : Culture et responsabilités}

Le point de départ du DevSecOps est organisationnel. La sécurité ne doit pas
être portée uniquement par une équipe spécialisée. Les développeurs,
administrateurs système, ingénieurs plateforme, responsables sécurité et
exploitants doivent partager une même responsabilité.

\begin{longtable}{|p{4cm}|p{9cm}|}
\hline
\textbf{Rôle} & \textbf{Responsabilité DevSecOps} \\
\hline
Développeur & Écrire du code testable, éviter les secrets, corriger les alertes
de sécurité proches du code. \\
\hline
DevOps / Platform Engineer & Automatiser les pipelines, standardiser les
environnements, fiabiliser les déploiements. \\
\hline
Security Engineer & Définir les contrôles, prioriser les risques, valider les
politiques de sécurité. \\
\hline
SRE & Mesurer la disponibilité, superviser les incidents, vérifier la résilience
et les objectifs de service. \\
\hline
\end{longtable}

\section{Étape 2 : Gestion du code source}

La première base technique est un dépôt propre, versionné et contrôlé. Chaque
modification doit passer par une revue et un pipeline automatisé.

\begin{itemize}[leftmargin=*]
  \item utiliser Git comme source de vérité ;
  \item protéger la branche principale ;
  \item exiger une revue avant fusion ;
  \item éviter les commits directs en production ;
  \item ajouter un fichier \texttt{.gitignore} adapté ;
  \item refuser les secrets, fichiers temporaires et artefacts sensibles.
\end{itemize}

\begin{lstlisting}[style=shellstyle,caption={Commandes Git de base}]
git status
git checkout -b feature/devsecops-bootstrap
git add .
git commit -m "Add DevSecOps bootstrap checks"
git push origin feature/devsecops-bootstrap
\end{lstlisting}

\section{Étape 3 : Intégration continue}

L'intégration continue permet de détecter rapidement les erreurs. Un pipeline
minimal doit exécuter les tests, vérifier la syntaxe, contrôler la structure du
projet et produire un résultat clair.

\begin{longtable}{|p{5cm}|p{8cm}|}
\hline
\textbf{Contrôle} & \textbf{But} \\
\hline
Lint & Détecter les erreurs de style, syntaxe et scripts invalides. \\
\hline
Tests unitaires & Vérifier les fonctions isolées du code applicatif. \\
\hline
Tests d'intégration & Vérifier les interactions entre composants. \\
\hline
Build conteneur & Vérifier que l'application peut être empaquetée. \\
\hline
Archivage des résultats & Garder les preuves du pipeline. \\
\hline
\end{longtable}

\section{Étape 4 : Sécurité shift-left}

Le principe \textit{shift-left} consiste à déplacer les contrôles de sécurité le
plus tôt possible dans le cycle de développement. Cela réduit le coût de
correction et évite de découvrir les problèmes au moment du déploiement.

\begin{itemize}[leftmargin=*]
  \item scan des dépendances ;
  \item scan des images Docker ;
  \item détection de secrets ;
  \item analyse statique de code ;
  \item validation des fichiers Kubernetes ;
  \item contrôle des configurations dangereuses.
\end{itemize}

\begin{lstlisting}[style=shellstyle,caption={Exemples de contrôles sécurité}]
trivy fs .
trivy image localhost:5001/mon-application:dev
syft localhost:5001/mon-application:dev -o cyclonedx-json
\end{lstlisting}

\section{Étape 5 : Conteneurisation sécurisée}

Les images conteneur doivent être reproductibles, minimales et durcies. Une
image de production ne doit pas contenir d'outils inutiles, de secrets ou de
dépendances de développement.

\begin{itemize}[leftmargin=*]
  \item utiliser une image de base maintenue ;
  \item exécuter l'application avec un utilisateur non root ;
  \item limiter les permissions du conteneur ;
  \item séparer dépendances de build et dépendances de runtime ;
  \item scanner l'image avant promotion ;
  \item éviter les tags mutables pour les releases finales.
\end{itemize}

\section{Étape 6 : Premiers contrôles Kubernetes}

Dans Kubernetes, la sécurité doit être décrite sous forme de manifestes et de
politiques. Les premiers contrôles à mettre en place concernent les ressources,
les permissions, les probes et l'exposition réseau.

\begin{longtable}{|p{5cm}|p{8cm}|}
\hline
\textbf{Contrôle Kubernetes} & \textbf{Intérêt} \\
\hline
Requests et limits & Éviter qu'un pod consomme toutes les ressources. \\
\hline
Readiness et liveness probes & Détecter les pods indisponibles. \\
\hline
SecurityContext & Interdire les privilèges inutiles. \\
\hline
NetworkPolicy & Réduire les communications réseau non nécessaires. \\
\hline
RBAC minimal & Limiter les permissions accordées aux services. \\
\hline
PodDisruptionBudget & Maintenir une disponibilité minimale pendant les
opérations. \\
\hline
\end{longtable}

\section{Étape 7 : Gestion des secrets}

Un secret ne doit jamais être stocké en clair dans le dépôt Git. La première
règle est de séparer les modèles de configuration des valeurs réelles.

\begin{itemize}[leftmargin=*]
  \item versionner uniquement des fichiers \texttt{.template} ou exemples ;
  \item stocker les vraies valeurs dans un gestionnaire de secrets ;
  \item limiter les permissions des fichiers locaux ;
  \item renouveler régulièrement les secrets ;
  \item documenter la procédure de rotation.
\end{itemize}

\section{Étape 8 : Supply chain logicielle}

Une chaîne DevSecOps crédible doit prouver l'origine et l'intégrité des
artefacts produits. Les étapes essentielles sont les suivantes :

\begin{enumerate}[leftmargin=*]
  \item construire l'image ;
  \item scanner l'image avec Trivy ;
  \item générer un SBOM avec Syft ;
  \item signer l'image avec Cosign ;
  \item vérifier la signature ;
  \item promouvoir l'image par digest immuable ;
  \item générer une attestation de release ;
  \item générer une provenance.
\end{enumerate}

\begin{lstlisting}[style=shellstyle,caption={Exemple de chaîne supply chain}]
make supply-chain-execute
make release-proof-strict
\end{lstlisting}

\section{Étape 9 : Observabilité et preuves}

Un système DevSecOps ne doit pas seulement fonctionner : il doit pouvoir le
prouver. Les métriques, journaux, rapports et snapshots permettent de défendre
l'état réel du système.

\begin{itemize}[leftmargin=*]
  \item état des nœuds Kubernetes ;
  \item état des pods et deployments ;
  \item métriques CPU et mémoire ;
  \item état HPA ;
  \item rapports Kyverno ;
  \item rapports Trivy, Syft et Cosign ;
  \item attestation et provenance de release ;
  \item support pack final.
\end{itemize}

\section{Étape 10 : Lecture honnête des statuts}

Pour éviter les fausses déclarations de réussite, chaque preuve doit être
classée avec un statut explicite.

\begin{longtable}{|p{4cm}|p{9cm}|}
\hline
\textbf{Statut} & \textbf{Signification} \\
\hline
TERMINÉ & La preuve réelle existe et le contrôle a réussi. \\
\hline
PARTIEL & Le contrôle a été exécuté mais un écart reste présent. \\
\hline
PRÊT\_NON\_EXÉCUTÉ & Le composant est prêt, mais l'action n'a pas encore été
lancée. \\
\hline
DÉPENDANT\_DE\_L\_ENVIRONNEMENT & Le résultat dépend d'un service externe, d'un
outil ou d'une infrastructure absente. \\
\hline
\end{longtable}

\section{Conclusion}

Les premiers pas DevSecOps consistent à rendre le projet testable, observable,
sécurisé et vérifiable dès le début. Une fois ces bases en place, il devient
possible de progresser vers une plateforme plus crédible : déploiement par
digest immuable, politiques Kubernetes, attestations de release, provenance,
haute disponibilité et résilience des données.

La valeur principale du DevSecOps est la preuve. Un système n'est pas seulement
déclaré sécurisé ou prêt pour la production : il doit produire des artefacts
concrets permettant de le démontrer.

\end{document}
